%-------------------------
% Resume in Latex
% Author : Jake Gutierrez
% Based off of: https://github.com/sb2nov/resume
% License : MIT
%------------------------

\documentclass[letterpaper,11pt]{article}

\usepackage[utf8]{inputenc}
%\usepackage[T1]{fontenc}
\usepackage{latexsym}
\usepackage[empty]{fullpage}
\usepackage{titlesec}
\usepackage{marvosym}
\usepackage[usenames,dvipsnames]{color}
\usepackage{verbatim}
\usepackage{enumitem}
\usepackage[hidelinks]{hyperref}
\usepackage{fancyhdr}
\usepackage[brazil]{babel}
\usepackage{tabularx}
\usepackage{fontawesome5}
\usepackage{multicol}
\setlength{\multicolsep}{-3.0pt}
\setlength{\columnsep}{-1pt}
\input{glyphtounicode}


%----------FONT OPTIONS----------
% sans-serif
% \usepackage[sfdefault]{FiraSans}
% \usepackage[sfdefault]{roboto}
% \usepackage[sfdefault]{noto-sans}
% \usepackage[default]{sourcesanspro}

% serif
% \usepackage{CormorantGaramond}
% \usepackage{charter}


\pagestyle{fancy}
\fancyhf{} % clear all header and footer fields
\fancyfoot{}
\renewcommand{\headrulewidth}{0pt}
\renewcommand{\footrulewidth}{0pt}


% Adjust margins
\addtolength{\oddsidemargin}{-0.6in}
\addtolength{\evensidemargin}{-0.5in}
\addtolength{\textwidth}{1.19in}
\addtolength{\topmargin}{-.7in}
\addtolength{\textheight}{1.4in}

\urlstyle{same}

\raggedbottom
\raggedright
\setlength{\tabcolsep}{0in}

% Sections formatting
\titleformat{\section}{
  \vspace{-4pt}\scshape\raggedright\large\bfseries
}{}{0em}{}[\color{black}\titlerule \vspace{-5pt}]

% Ensure that generate pdf is machine readable/ATS parsable
\pdfgentounicode=1

%-------------------------
% Custom commands
\newcommand{\cert}[2]{\href{#1}{\scriptsize [certificado]}\ \href{#1}{#2}}
\newcommand{\certitem}[2]{\item\small \cert{#1}{#2}}

\newcommand{\resumeItem}[1]{
  \item\small{
    {#1 \vspace{-2pt}}
  }
}

\newcommand{\classesList}[4]{
    \item\small{
        {#1 #2 #3 #4 \vspace{-2pt}}
  }
}

\newcommand{\resumeSubheading}[4]{
  \vspace{-2pt}\item
    \begin{tabular*}{1.0\textwidth}[t]{l@{\extracolsep{\fill}}r}
      \textbf{#1} & \textbf{\small #2} \\
      \textit{\small#3} & \textit{\small #4} \\
    \end{tabular*}\vspace{-7pt}
}

\newcommand{\resumeSubSubheading}[2]{
    \item
    \begin{tabular*}{0.97\textwidth}{l@{\extracolsep{\fill}}r}
      \textit{\small#1} & \textit{\small #2} \\
    \end{tabular*}\vspace{-7pt}
}

\newcommand{\resumeProjectHeading}[2]{
    \item
    \begin{tabular*}{1.001\textwidth}{l@{\extracolsep{\fill}}r}
      \small#1 & \textbf{\small #2}\\
    \end{tabular*}\vspace{-7pt}
}

\newcommand{\resumeSubItem}[1]{\resumeItem{#1}\vspace{-4pt}}

\renewcommand\labelitemi{$\vcenter{\hbox{\tiny$\bullet$}}$}
\renewcommand\labelitemii{$\vcenter{\hbox{\tiny$\bullet$}}$}

\newcommand{\resumeSubHeadingListStart}{\begin{itemize}[leftmargin=0.0in, label={}]}
\newcommand{\resumeSubHeadingListEnd}{\end{itemize}}
\newcommand{\resumeItemListStart}{\begin{itemize}}
\newcommand{\resumeItemListEnd}{\end{itemize}\vspace{-5pt}}

%-------------------------------------------
%%%%%%  RESUME STARTS HERE  %%%%%%%%%%%%%%%%%%%%%%%%%%%%


\begin{document}

%----------HEADING----------
\begin{center}
    {\Huge \scshape Victor T. Rodrigues} \\ \vspace{1pt}
    Presidente Prudente, SP \\ \vspace{1pt}
 \small \raisebox{-0.1\height} \faPhone \ (18) 99681-1971~\href{mailto:victortavi@outlook.com}{\raisebox{-0.2\height}\faEnvelope\  \underline{victortavi@outlook.com}} ~
 \href{https://linkedin.com/in/victor-taveira}{\raisebox{-0.2\height}\faLinkedin\ \underline{linkedin.com/in/victor-taveira}}~
\href{https://github.com/elvitin}{\raisebox{-0.2\height}\faGithub\ \underline{github.com/elvitin}} ~
    %\href{https://taveira.dev}{\raisebox{-0.2\height}\faGlobe\ \underline{taveira.dev}}
    \vspace{-8pt}
\end{center}


%-----------EXPERIENCE-----------
\section{Experiência}
  \resumeSubHeadingListStart

    \resumeSubheading
      {Nimbus Sistemas Ltda}{Ago. 2025 -- Atual}
      {Lead Software Engineer}{Presidente Prudente, SP -- Remoto}
      \resumeItemListStart
        \resumeItem{Condução de entrevistas com stakeholders e estudo de ferramentas case de mercado para elicitação de requisitos.}
        \resumeItem{Criação de Repositório de Conhecimento (Wiki) formalizando a gestão do conhecimento do produto}
        \resumeItem{Aplicação de algumas cerimônias Scrum, criação de cronograma técnico}
        \resumeItem{Desenvolvimento de aplicação Monorepo - Software Jurídico que da apoio a criação de Laudo e execução de Perícias Judiciais}
        \resumeItem{Implementação de infraestrutura CI/CD com entregas automáticas na Azure utilizando GitHub Actions.}
        \resumeItemListEnd

    \resumeSubheading
      {Rubcube Tratamento de Dados Ltda}{Março 2024 -- Jun. 2025}
      {Engenheiro de Software Junior}{Presidente Prudente, SP}
      \resumeItemListStart
        %\resumeItem{Participei do Programa de Aceleração \href{https://rubcube.com.br/rubcamp}{\textcolor{blue}{\underline{Rubcamp}}}}
        %\resumeItem{Participei do Programa de Aceleração \href{https://rubcube.com/rubcamp/}{\textcolor{blue}{Rubcamp}}}
        \resumeItem{Participação no programa de aceleração \href{https://rubcube.com/rubcamp/}{\underline{Rubcamp}} onde obtive um dos melhores desempenhos de entrega de software}
        \resumeItem{Apoio técnico a participantes do \href{https://rubcube.com/rubcamp/}{\underline{Rubcamp}} em desafios de infraestrutura, conteinerização Docker, Deploy (CI/CD), Backend (NodeJs) e Mobile (React Native), contribuindo diretamente para seu avanço e aprovação no programa}
        \resumeItem{Atuação no time de desenvolvimento principal, Core Projetos em ReactJs, React Native, Backend NodeJs.}
        \resumeItem{Atuação na implementação de features no app mobile \href{https://play.google.com/store/apps/details?id=br.com.road.app}{\underline{Road}}.}
        \resumeItem{Atuação na implementação de features \& correções de bugs no app mobile de Benefícios \href{https://play.google.com/store/apps/details?id=net.natura.pay&hl=pt_BR}{\underline{Emanapay}} da Natura.}
        \resumeItem{Implementação de funcionalidades de Bloqueio de Cartão \& SMS Token no aplicativo mobile \href{https://www.bancodigimais.com.br/}{\underline{Digi+}}.}
        \resumeItem{Construção de consumidores e integração com filas AWS SQS.}
        \resumeItem{Construção de repositórios DynamoDB}
        \resumeItem{Implantação de Recursos de Infraestrutura via Terraforms}
        \resumeItem{Contato com ferramentas de infraestrutura, AWS API Gateway, AWS SSM, AWS Serverless, Google Firebase, AWS CDK, AWS CloudFormation, Terraform, Docker, Docker Compose, GitHub Actions.}
        \resumeItem{Automação com Shell Script, Python e JavaScript}
        \resumeItem{Utilização de qualidade de código: Testes Automatizados, (Integração, unidade, E2E), Ferramentas de testes Jest, Vitest e Cypress; Aplicação de Arquitetura Hexagonal e Clean Architecture.}
      \resumeItemListEnd

    \resumeSubheading
      {Prefeitura Municipal}{2018 -- 2020}
      {Estágio em Suporte Técnico -- Dep. de Educação}{Nantes, SP -- Híbrido}
      \resumeItemListStart
        \resumeItem{Execução de atividades Administrativas}
        \resumeItem{Prestação de Suporte Técnico}
        \resumeItem{Treinamento de Usuários no uso de ferramentas Administrativas}
      \resumeItemListEnd

    \resumeSubheading
      {Mercado Conceição}{2015 -- 2018}
      {Consultor de Implantação}{Nantes, SP -- Presencial}
      \resumeItemListStart
        \resumeItem{Levantamento de Requisitos de Infraestrutura para Automação do Mercado}
        \resumeItem{Criação de Orçamento de Infraestrutura, Compra de Equipamentos \& Software}
        \resumeItem{Instalação de Equipamentos, Leitores, Balanças, Homologação, Instalação de Software \& Parametrização de Software PDV/ERP}
        \resumeItem{Treinamento de Usuários/Colaboradores na Operação do Sistema}
        \resumeItem{Participação direta no aumento de resultados financeiros}
      \resumeItemListEnd

  \resumeSubHeadingListEnd
\vspace{-16pt}

%-----------Projetos-----------
\section{Projetos}
    \vspace{-5pt}
    \resumeSubHeadingListStart
    
      \resumeProjectHeading
          {\textbf{\href{https://github.com/elvitin/devtime-releases}{\underline{DevTime}}} $|$ \emph{Rust, ReactJs, Typescript, \href{https://tauri.app/}{\underline{Tauri 2.0}}, \href{https://ui.shadcn.com/}{\underline{Shadcn}}}}{Fevereiro 2026}
          \resumeItemListStart
            \resumeItem{Ferramenta Desktop feita com ReactJs e Rust usando Tauri}
            \resumeItem{Surgiu devia a uma necessidade emergente onde atuando como Pessoa Jurídica, houve-se a necessidade de registrar a carga horária com mais precisão e facilidade}
            \resumeItem{Antes, o registros de tempo gasto em tarefas de desenvolvimento eram registrados em uma planilha de forma manual}
            \resumeItem{Com a ferramenta, é possível contabilizar todo o horário produtivo, permitindo pausas e retomadas}
            \resumeItem{Além disso, a ferramenta exporta uma planilha já com os custos das horas de desenvolvimento acumuladas com precisão}
            \resumeItem{A entrega do app ocorre via CI/CD através de um pipeline de build e distribuição no Github Actions, onde toda nova versão é publicada em \href{https://github.com/elvitin/devtime-releases}{\underline{devtime-releases}}}
          \resumeItemListEnd
          \vspace{-13pt}
          
      \resumeProjectHeading
          {\textbf{\href{https://github.com/elvitin/ezychat}{\underline{Ezychat}}} $|$ \emph{Java, JavaFx, Maven, Protobuf, ORMLite, Slf4j}}{Julho 2025}
          \resumeItemListStart
            \resumeItem{Projeto que nasceu a partir de conhecimentos obtidos da disciplina de Sistemas Operacionais \& Redes de Computadores}
            \resumeItem{Esse projeto explora conceitos base para entendimento de protocolos em diferentes camadas de rede, gerenciamento de processos e threads, etc.}
            \resumeItem{É uma aplicação de mensagens instantâneas em tempo real que imita no WhatsApp, com arquitetura cliente-servidor, utilizando Java 24, JavaFX para a interface gráfica e comunicação via TCP Sockets com serialização Protocol Buffers (Protobuf), garantindo alta performance e baixa latência na troca de mensagens.}
            \resumeItem{Estruturei o projeto com arquitetura modular multi-módulo Maven (server, app, common), aplicando princípios de Clean Architecture com separação clara entre camadas de domínio, infraestrutura, casos de uso e apresentação, facilitando manutenção e escalabilidade.}
            \resumeItem{Implementei comunicação assíncrona e concorrente utilizando Virtual Threads (Project Loom), ConcurrentHashMap, BlockingQueue e CopyOnWriteArrayList para gerenciamento eficiente de conexões simultâneas e processamento de mensagens sem bloqueio.}
            \resumeItem{Desenvolvi sistema completo de autenticação e registro de usuários com validações de negócio (unicidade de username/email, requisitos de senha), notificação de presença online em tempo real e restauração automática de mensagens pendentes ao reconectar.}
            \resumeItem{Implementei padrão Event Bus (publish/subscribe) para desacoplamento entre componentes da aplicação cliente, permitindo comunicação reativa entre a camada de rede e a interface gráfica sem acoplamento direto.}
            \resumeItem{Integrei armazenamento de arquivos com AWS S3 para upload de fotos de perfil dos usuários, e SQLite com ORM (ORMLite) para persistência local de dados no servidor.}
            \resumeItem{Defini contrato de comunicação entre cliente e servidor com Protocol Buffers, criando schemas .proto tipados para login, registro, mensagens privadas, status de usuário e confirmações — garantindo interoperabilidade e versionamento seguro do protocolo.}
            \resumeItem{Construí interface gráfica moderna com JavaFX + FXML + CSS, incluindo componentes customizados como lista de mensagens, perfis de contatos com imagens arredondadas e chat aberto, seguindo padrão MVC com controllers dedicados}
            \resumeItem{Apliquei design patterns como Factory (criação de use cases e repositórios), Singleton (gerenciador de conexões e event bus), Strategy (handlers de payload) e Gateway (abstração de mensageria), demonstrando sólido conhecimento em engenharia de software.}
            \resumeItem{Tecnologias-chave para destacar: Java 24, Virtual Threads, Protocol Buffers, TCP Sockets, JavaFX, Maven Multi-Module, Clean Architecture, AWS S3, SQLite, ORMLite, Design Patterns, Programação Concorrente.}
          \resumeItemListEnd
          \vspace{-13pt}
        
      \resumeProjectHeading
          {\textbf{\href{https://github.com/elvitin/rubbank-backend}{\underline{RubBank API REST de Banco Digital (Pré IA)}}} $|$ \emph{NodeJs, Typescript, MVC, JWT}}{Março 2024}
          \resumeItemListStart
            \resumeItem{RubBank Backend é uma API REST de banco digital com Node.js, TypeScript, Express, Prisma/PostgreSQL, autenticação JWT com invalidação dinâmica de tokens, criptografia bcrypt, validação com Zod e cobertura de testes com Vitest (unitários, integração e E2E). Ambiente dockerizado com documentação OpenAPI/Swagger.}
            \resumeItem{Surgiu durante a participação do Rubcamp, esse backend é onde concentra-se a lógica de negócios necessárias para o app \href{https://github.com/elvitin/rubbank-mobile/}{\underline{RubBank Mobile}} }
            \resumeItem{Ambiente 100\% containerizado com Docker, facilitando reprodutibilidade entre máquinas.}
            \resumeItem{Testes em três camadas: unitários (validadores), integração (banco de dados) e E2E (endpoints HTTP reais).}
            \resumeItem{JWT com segredo dinâmico: o segredo é derivado de chave\_base + hash\_senha + timestamp\_último\_logout, tornando tokens inválidos automaticamente após logout ou troca de senha, sem necessidade de lista negra.}
            \resumeItem{Transações atômicas com Prisma (\$transaction) para garantir consistência nas transferências financeiras.}
          \resumeItemListEnd
          \vspace{-13pt}
            
        \resumeProjectHeading
    {\textbf{\href{https://github.com/elvitin/zupzap}{\underline{ZupZap}}} $|$ \emph{C\#, ASP.NET Core, React, MySQL, Docker Compose, Make}}{Dezembro 2022}
    \resumeItemListStart
      \resumeItem{Desenvolvi uma aplicação full stack de mensagens com frontend em React 18 e backend em ASP.NET Core, expondo API REST para cadastro, listagem, consulta e exclusão de mensagens.}
      \resumeItem{Estruturei o backend em camadas Controller, Service e Repository, separando regras de negócio, validações e acesso a dados em MySQL.}
      \resumeItem{Implementei persistência relacional para mensagens e script de inicialização do banco, com modelo contendo identificador, texto, usuário e data de envio.}
      \resumeItem{Integrei frontend e backend por meio de proxy para rotas /api e cliente HTTP reutilizável, simplificando o consumo dos endpoints no ambiente local e conteinerizado.}
      \resumeItem{Containerizei banco, API e frontend com Docker Compose, utilizando healthcheck e automação via Makefile para subir toda a stack e validar a prontidão dos serviços com um comando.}
      \resumeItem{Apliquei validações de entrada e tratamento básico de conteúdo no serviço de domínio antes da gravação das mensagens.}
    \resumeItemListEnd
    \vspace{-13pt}
    
    \resumeSubHeadingListEnd
\vspace{-16pt}


%-----------EDUCATION-----------
\section{Educação}
  \resumeSubHeadingListStart
    \resumeSubheading
      {FIPP -- Universidade do Oeste Paulista}{Fev. 2022 -- Atual}
      {Sistemas de Informação -- Bacharelado}{Presidente Prudente, SP}
        \resumeItemListStart
            \resumeItem{Formação com foco em Engenharia de Software, Desenvolvimento Web/Mobile e Arquitetura de Sistemas.}
            \resumeItem{Base em Banco de Dados, Estruturas de Dados, Algoritmos, Pesquisa e Ordenação e Pesquisa Operacional.}
            \resumeItem{Conhecimentos em Redes de Computadores, Sistemas Operacionais, Segurança e Auditoria de Sistemas.}
            \resumeItem{Experiência acadêmica em IA, BI, MLOps, Computação em Nuvem e integração de serviços.}
            \resumeItem{Disciplinas de gestão: Gestão de Projetos, Governança de TI, Operações e Logística, Simulação Gerencial.}
            \resumeItem{Projetos práticos e Estágio Supervisionado com aplicação de conceitos em cenários reais.}
        \resumeItemListEnd
    \resumeSubheading
      {Escola Estadual Rage Andeiras Prof.}{Concluído}
      {Ensino Médio}{Nantes, SP}
  \resumeSubHeadingListEnd

%-----------TECHNICAL SKILLS-----------
% \section{Habilidades Técnicas}
%  \begin{itemize}[leftmargin=0.15in, label={}]
%     \small{\item{
%      \textbf{Linguagens}{: JavaScript, Java, C\#, C, C++, SQL, Python, Bash} \\
%      \textbf{Ferramentas}{: Git, GitHub Actions, Docker, Docker Compose, Figma, Adobe XD, Firebase} \\
%      \textbf{Tecnologias/Frameworks}{: Node.js, Deno, Express.js, NestJS, Java EE, Spring, Spring Boot, React Native, ReactJS, HTML5, CSS3, Bootstrap, Tailwind CSS, PostgreSQL, Oracle, Linux, AWS, RESTful APIs, Jest, Vitest, Cypress} \\
%     }}
%  \end{itemize}
%  \vspace{-16pt}

%-----------QUALIFICAÇÕES-----------
% \section{Qualificações}
%     \resumeSubHeadingListStart
%         \resumeSubheading{Desenvolvimento de Software}{Resumo}{Competências principais}{ }
%             \resumeItemListStart
%                 \resumeItem{Experiência inicial no desenvolvimento de APIs com Node.js e Express.js, além de Java EE com Servlets e Spring Boot.}
%                 \resumeItem{Conhecimento em Linux, arquitetura RESTful, lógica de programação e noções de arquitetura de computadores.}
%                 \resumeItem{Perfil com proatividade, boa comunicação, capacidade de resolução de problemas e adaptação a novas tecnologias.}
%             \resumeItemListEnd
%     \resumeSubHeadingListEnd


%------CURSOS COMPLEMENTARES-------
\section{Formação Complementar \& Certificações}
\begin{itemize}[leftmargin=0.15in, label={}]

    \item \small \textbf{\href{https://anthropic.skilljar.com/}{\underline{Anthropic Skilljar}}}
    \begin{itemize}[itemsep=-5pt, parsep=3pt]
        \item\small \href{https://verify.skilljar.com/c/26xevfmxt6xn}{\underline{Claude Code in Action -- Março 2026}}
    \end{itemize}


    \item \small \textbf{\href{https://raw.githubusercontent.com/elvitin/curriculo/main/certs/infoeste/Certificado2024.pdf}{\underline{Infoeste 2024 -- 36ª Semana de Computação e Informática da FIPP/Unoeste}}}
    \begin{itemize}[itemsep=-5pt, parsep=3pt]
        \item\small Global Career Opportunities - Construindo uma Carreira Internacional: Desbloqueie seu Potencial para uma Carreira em TI de Sucesso!
        \item\small Google Cloud - Crie Modelos de Machine Learning com o BigQuery ML
        \item\small Arduino e Tinkercad: Dominando a Eletrônica Simulada
        \item\small Produção de conteúdo para o Instagram
    \end{itemize}

    \item \small \textbf{\href{https://raw.githubusercontent.com/elvitin/curriculo/main/certs/infoeste/Certificado2023.pdf}{\underline{Infoeste 2023 -- 35ª Semana de Computação e Informática da FIPP/Unoeste}}}
    \begin{itemize}[itemsep=-5pt, parsep=3pt]
        \item\small Google - Conceitos básicos de  computação em nuvem
        \item\small Google Cloud - Crie Modelos de Machine Learning com o BigQuery ML
        \item\small Arduino e Tinkercad: Dominando a Eletrônica Simulada
        \item\small Produção de conteúdo para o Instagram
    \end{itemize}

    %     \item \small \textbf{\href{https://raw.githubusercontent.com/elvitin/curriculo/main/certs/desenvolve/CertificadoDesenvolve2023.pdf}{\underline{Programa Desenvolve 2023 - Formação em FullStack}}}
    % \begin{itemize}[itemsep=-5pt, parsep=3pt]
    %     \item\small \href{https://cursos.alura.com.br/user/v-rodrigues/course/oracle-database-sql-dml/formalCertificate}{SQL com Oracle: manipule e consulte dados}
    %     \item\small \href{https://cursos.alura.com.br/user/v-rodrigues/course/mysql-manipule-dados-com-sql/formalCertificate}{SQL com MySQL: manipule e consulte dados}
    %     \item\small \href{https://cursos.alura.com.br/user/v-rodrigues/course/shellscripting/formalCertificate}{Shell Scripting parte 1: scripts de automação de tarefas}
    %     \item\small \href{https://cursos.alura.com.br/user/v-rodrigues/course/node-rest-api/formalCertificate}{Rest com NodeJS: API com Express e MySQL}
    %     \item\small \href{https://cursos.alura.com.br/user/v-rodrigues/course/introducao-postgresql-primeiros-passos/formalCertificate}{PostgreSQL}
    %     \item\small \href{https://cursos.alura.com.br/user/v-rodrigues/course/orm-nodejs-avancando-sequelize/formalCertificate}{ORM com NodeJS: avançando nas funcionalidades do Sequelize}
    %     \item\small \href{https://cursos.alura.com.br/user/v-rodrigues/course/orm-nodejs-api-sequelize-mysql/formalCertificate}{ORM com NodeJS: API com Sequelize e MySQL}
    %     \item\small \href{https://cursos.alura.com.br/user/v-rodrigues/course/oracle-plsql-procedures-funcoes-excecoes/formalCertificate}{Oracle PL/SQL: Procedures, funções e exceções}
    %     \item\small \href{https://cursos.alura.com.br/user/v-rodrigues/course/nodejs-streaming-dados/formalCertificate}{NodeJS: Streaming de dados e Repositório}
    %     \item\small \href{https://cursos.alura.com.br/user/v-rodrigues/course/nodejs-api-rest-padronizada-escalavel/formalCertificate}{NodeJS: crie uma API REST padronizada e escalável}
    %     \item\small \href{https://cursos.alura.com.br/user/v-rodrigues/course/nodejs-api-rest-controle-versao/formalCertificate}{NodeJS: avançando em APIs REST com controle de versões}
    %     \item\small \href{https://cursos.alura.com.br/user/v-rodrigues/course/maven-gerenciamento-dependencias-build-aplicacoes-java/formalCertificate}{Maven: gerenciamento de dependências e build de aplicações Java}
    %     \item\small \href{https://cursos.alura.com.br/user/v-rodrigues/course/fundamentos-javascript-tipos-variaveis-funcoes/formalCertificate}{JavaScript: tipos, variáveis e funções}
    %     \item\small \href{https://cursos.alura.com.br/user/v-rodrigues/course/fundamentos-javascript-objetos/formalCertificate}{JavaScript: objetos}
    %     \item\small \href{https://cursos.alura.com.br/user/v-rodrigues/course/servlets-fundamentos-programacao-web-java/formalCertificate}{Java Servlet: programação web Java}
    %     \item\small \href{https://cursos.alura.com.br/user/v-rodrigues/course/servlet-autenticacao-autorizacao-mvc/formalCertificate}{Java Servlet: autenticação, autorização e o padrão MVC}
    %     \item\small \href{https://cursos.alura.com.br/user/v-rodrigues/course/http-fundamentos/formalCertificate}{HTTP: Entendendo a web por baixo dos panos}
    %     \item\small \href{https://cursos.alura.com.br/user/v-rodrigues/course/git-github-controle-de-versao/formalCertificate}{Git e Github: controle e compartilhe seu código}
    %     \item\small \href{https://cursos.alura.com.br/user/v-rodrigues/course/habitos/formalCertificate}{Hábitos: da produtividade às metas pessoais}
    %     \item\small \href{https://cursos.alura.com.br/user/v-rodrigues/course/feedback/formalCertificate}{Feedback: a arte de orientar e ser orientado com sentido}
    %     \item\small \href{https://cursos.alura.com.br/user/v-rodrigues/course/comunicacao/formalCertificate}{Comunicação: como se expressar bem e ser compreendido}
    % \end{itemize}

%     \item \small \textbf{\href{https://raw.githubusercontent.com/elvitin/curriculo/main/certs/desenvolve/CertificadoDesenvolve2023.pdf}{\underline{Programa Desenvolve 2023 - Formação em FullStack}}}
% \begin{itemize}[itemsep=-5pt, parsep=3pt]
%     \item\small \href{https://cursos.alura.com.br/user/v-rodrigues/course/oracle-database-sql-dml/formalCertificate}{\underline{SQL com Oracle: manipule e consulte dados}}
%     \item\small \href{https://cursos.alura.com.br/user/v-rodrigues/course/mysql-manipule-dados-com-sql/formalCertificate}{\underline{SQL com MySQL: manipule e consulte dados}}
%     \item\small \href{https://cursos.alura.com.br/user/v-rodrigues/course/shellscripting/formalCertificate}{\underline{Shell Scripting parte 1: scripts de automação de tarefas}}
%     \item\small \href{https://cursos.alura.com.br/user/v-rodrigues/course/node-rest-api/formalCertificate}{\underline{Rest com NodeJS: API com Express e MySQL}}
%     \item\small \href{https://cursos.alura.com.br/user/v-rodrigues/course/introducao-postgresql-primeiros-passos/formalCertificate}{\underline{PostgreSQL}}
%     \item\small \href{https://cursos.alura.com.br/user/v-rodrigues/course/orm-nodejs-avancando-sequelize/formalCertificate}{\underline{ORM com NodeJS: avançando nas funcionalidades do Sequelize}}
%     \item\small \href{https://cursos.alura.com.br/user/v-rodrigues/course/orm-nodejs-api-sequelize-mysql/formalCertificate}{\underline{ORM com NodeJS: API com Sequelize e MySQL}}
%     \item\small \href{https://cursos.alura.com.br/user/v-rodrigues/course/oracle-plsql-procedures-funcoes-excecoes/formalCertificate}{\underline{Oracle PL/SQL: Procedures, funções e exceções}}
%     \item\small \href{https://cursos.alura.com.br/user/v-rodrigues/course/nodejs-streaming-dados/formalCertificate}{\underline{NodeJS: Streaming de dados e Repositório}}
%     \item\small \href{https://cursos.alura.com.br/user/v-rodrigues/course/nodejs-api-rest-padronizada-escalavel/formalCertificate}{\underline{NodeJS: crie uma API REST padronizada e escalável}}
%     \item\small \href{https://cursos.alura.com.br/user/v-rodrigues/course/nodejs-api-rest-controle-versao/formalCertificate}{\underline{NodeJS: avançando em APIs REST com controle de versões}}
%     \item\small \href{https://cursos.alura.com.br/user/v-rodrigues/course/maven-gerenciamento-dependencias-build-aplicacoes-java/formalCertificate}{\underline{Maven: gerenciamento de dependências e build de aplicações Java}}
%     \item\small \href{https://cursos.alura.com.br/user/v-rodrigues/course/fundamentos-javascript-tipos-variaveis-funcoes/formalCertificate}{\underline{JavaScript: tipos, variáveis e funções}}
%     \item\small \href{https://cursos.alura.com.br/user/v-rodrigues/course/fundamentos-javascript-objetos/formalCertificate}{\underline{JavaScript: objetos}}
%     \item\small \href{https://cursos.alura.com.br/user/v-rodrigues/course/servlets-fundamentos-programacao-web-java/formalCertificate}{\underline{Java Servlet: programação web Java}}
%     \item\small \href{https://cursos.alura.com.br/user/v-rodrigues/course/servlet-autenticacao-autorizacao-mvc/formalCertificate}{\underline{Java Servlet: autenticação, autorização e o padrão MVC}}
%     \item\small \href{https://cursos.alura.com.br/user/v-rodrigues/course/http-fundamentos/formalCertificate}{\underline{HTTP: Entendendo a web por baixo dos panos}}
%     \item\small \href{https://cursos.alura.com.br/user/v-rodrigues/course/git-github-controle-de-versao/formalCertificate}{\underline{Git e Github: controle e compartilhe seu código}}
%     \item\small \href{https://cursos.alura.com.br/user/v-rodrigues/course/habitos/formalCertificate}{\underline{Hábitos: da produtividade às metas pessoais}}
%     \item\small \href{https://cursos.alura.com.br/user/v-rodrigues/course/feedback/formalCertificate}{\underline{Feedback: a arte de orientar e ser orientado com sentido}}
%     \item\small \href{https://cursos.alura.com.br/user/v-rodrigues/course/comunicacao/formalCertificate}{\underline{Comunicação: como se expressar bem e ser compreendido}}
% \end{itemize}

    \item \small \textbf{\href{https://raw.githubusercontent.com/elvitin/curriculo/main/certs/desenvolve/CertificadoDesenvolve2023.pdf}{\underline{Programa Desenvolve 2023 - Formação em FullStack}}}
\begin{itemize}[itemsep=-5pt, parsep=3pt]
    \item\small \cert{https://cursos.alura.com.br/user/v-rodrigues/course/oracle-database-sql-dml/formalCertificate}{SQL com Oracle: manipule e consulte dados}
    \item\small \cert{https://cursos.alura.com.br/user/v-rodrigues/course/mysql-manipule-dados-com-sql/formalCertificate}{SQL com MySQL: manipule e consulte dados}
    \item\small \cert{https://cursos.alura.com.br/user/v-rodrigues/course/shellscripting/formalCertificate}{Shell Scripting parte 1: scripts de automação de tarefas}
    \item\small \cert{https://cursos.alura.com.br/user/v-rodrigues/course/node-rest-api/formalCertificate}{Rest com NodeJS: API com Express e MySQL}
    \item\small \cert{https://cursos.alura.com.br/user/v-rodrigues/course/introducao-postgresql-primeiros-passos/formalCertificate}{PostgreSQL}
    \item\small \cert{https://cursos.alura.com.br/user/v-rodrigues/course/orm-nodejs-avancando-sequelize/formalCertificate}{ORM com NodeJS: avançando nas funcionalidades do Sequelize}
    \item\small \cert{https://cursos.alura.com.br/user/v-rodrigues/course/orm-nodejs-api-sequelize-mysql/formalCertificate}{ORM com NodeJS: API com Sequelize e MySQL}
    \item\small \cert{https://cursos.alura.com.br/user/v-rodrigues/course/oracle-plsql-procedures-funcoes-excecoes/formalCertificate}{Oracle PL/SQL: Procedures, funções e exceções}
    \item\small \cert{https://cursos.alura.com.br/user/v-rodrigues/course/nodejs-streaming-dados/formalCertificate}{NodeJS: Streaming de dados e Repositório}
    \item\small \cert{https://cursos.alura.com.br/user/v-rodrigues/course/nodejs-api-rest-padronizada-escalavel/formalCertificate}{NodeJS: crie uma API REST padronizada e escalável}
    \item\small \cert{https://cursos.alura.com.br/user/v-rodrigues/course/nodejs-api-rest-controle-versao/formalCertificate}{NodeJS: avançando em APIs REST com controle de versões}
    \item\small \cert{https://cursos.alura.com.br/user/v-rodrigues/course/maven-gerenciamento-dependencias-build-aplicacoes-java/formalCertificate}{Maven: gerenciamento de dependências e build de aplicações Java}
    \item\small \cert{https://cursos.alura.com.br/user/v-rodrigues/course/fundamentos-javascript-tipos-variaveis-funcoes/formalCertificate}{JavaScript: tipos, variáveis e funções}
    \item\small \cert{https://cursos.alura.com.br/user/v-rodrigues/course/fundamentos-javascript-objetos/formalCertificate}{JavaScript: objetos}
    \item\small \cert{https://cursos.alura.com.br/user/v-rodrigues/course/servlets-fundamentos-programacao-web-java/formalCertificate}{Java Servlet: programação web Java}
    \item\small \cert{https://cursos.alura.com.br/user/v-rodrigues/course/servlet-autenticacao-autorizacao-mvc/formalCertificate}{Java Servlet: autenticação, autorização e o padrão MVC}
    \item\small \cert{https://cursos.alura.com.br/user/v-rodrigues/course/http-fundamentos/formalCertificate}{HTTP: Entendendo a web por baixo dos panos}
    \item\small \cert{https://cursos.alura.com.br/user/v-rodrigues/course/git-github-controle-de-versao/formalCertificate}{Git e Github: controle e compartilhe seu código}
    \item\small \cert{https://cursos.alura.com.br/user/v-rodrigues/course/habitos/formalCertificate}{Hábitos: da produtividade às metas pessoais}
    \item\small \cert{https://cursos.alura.com.br/user/v-rodrigues/course/feedback/formalCertificate}{Feedback: a arte de orientar e ser orientado com sentido}
    \item\small \cert{https://cursos.alura.com.br/user/v-rodrigues/course/comunicacao/formalCertificate}{Comunicação: como se expressar bem e ser compreendido}
\end{itemize}

    \item \small \textbf{\href{https://raw.githubusercontent.com/elvitin/curriculo/main/certs/infoeste/Certificado2022.pdf}{\underline{Infoeste 2022 -- 34ª Semana de Computação e Informática da FIPP/Unoeste}}}
    \begin{itemize}[itemsep=-5pt, parsep=3pt]
        \item\small Empreendedorismo Tecnológico
        \item\small Marketing Digital
        \item\small Introdução ao Agile Thinking
        \item\small Power BI
    \end{itemize}
    \end{itemize}

%------ Outros -------
\section{Palestras, Webinars, Eventos \& Hackathons}
\begin{itemize}[leftmargin=0.15in, label={}]

    \item \small \textbf{\href{https://raw.githubusercontent.com/elvitin/curriculo/main/certs/sebrae-startupday/certificado_startupday_marco_2026.pdf}{\underline{Março 2026 - Startup Day - SEBRAE Startups}}}
    
    
    \item \small \textbf{\href{https://raw.githubusercontent.com/elvitin/curriculo/main/certs/infoeste/Certificado2025.pdf}{\underline{Infoeste 2025 -- 37ª Semana de Computação e Informática da FIPP/Unoeste}}}
    \begin{itemize}[itemsep=-5pt, parsep=3pt]
        \item\small (Hackathon) 19ª Maratona de Programação da FIPP/Unoeste
    \end{itemize}

    \item \small \textbf{\href{https://raw.githubusercontent.com/elvitin/curriculo/main/certs/infoeste/Certificado2024.pdf}{\underline{Infoeste 2024 -- 36ª Semana de Computação e Informática da FIPP/Unoeste}}}
    \begin{itemize}[itemsep=-5pt, parsep=3pt]
        \item\small (Hackathon) 18ª Maratona de Programação da FIPP/Unoeste
        \item\small (Evento) 14º Festival Cultural FIPP/Unoeste
        \item\small (Palestra) IA Generativa: Criando o Novo com Inteligência Artificial
        \item\small (Palestra) Cloud Run: From Zero to Hero (GDG DevFest)
    \end{itemize}

    \item \small \textbf{\href{https://raw.githubusercontent.com/elvitin/curriculo/main/certs/infoeste/Certificado2023.pdf}{\underline{Infoeste 2023 -- 35ª Semana de Computação e Informática da FIPP/Unoeste}}}
    \begin{itemize}[itemsep=-5pt, parsep=3pt]
        \item\small (Hackathon) 17ª Maratona de Programação da FIPP/Unoeste
        \item\small (Evento) 13º Festival Cultural FIPP/Unoeste
        \item\small (Palestra) Do Home Office para o Mundo!!!!
        \item\small (Palestra) IA Generativa com Large Language Models: uma introdução "mão-na-massa"
        \item\small (Palestra) Ipag: A trajetória de uma Fintech de Sucesso
        \item\small (Evento) Modelagem de Dados Dimensional para BI e Analytics (GDG DevFest)
    \end{itemize}

    \item \small \textbf{\href{https://raw.githubusercontent.com/elvitin/curriculo/main/certs/infoeste/Certificado2022.pdf}{\underline{Infoeste 2022 -- 34ª Semana de Computação e Informática da FIPP/Unoeste}}}
    \begin{itemize}[itemsep=-5pt, parsep=3pt]
        \item\small (Palestra) B3, Havan, Itaú, Magalu Smarthint e Vericode EXPERIENCE
        \item\small (Hackathon) 16ª Maratona de Programação da FIPP/Unoeste
    \end{itemize}

    \item \small \textbf{\href{https://raw.githubusercontent.com/elvitin/curriculo/main/certs/infoeste/Certificado2020.pdf}{\underline{Infoeste 2020 -- 33ª Semana de Computação e Informática da FIPP/Unoeste}}}
    \begin{itemize}[itemsep=-5pt, parsep=3pt]
        \item\small (Webinar) FIPP/Unoeste Mundo
        \item\small (Webinar) Forense computacional no âmbito da perícia oficial
        \item\small (Webinar) O profissional de Data Science
        \item\small (Webinar) Tornando-se um Gamedev
    \end{itemize}

    \item \small \textbf{\href{https://raw.githubusercontent.com/elvitin/curriculo/main/certs/infoeste/Certificado2019.pdf}{\underline{Infoeste 2019 -- 32ª Semana de Computação e Informática da FIPP/Unoeste}}}
    \begin{itemize}[itemsep=-5pt, parsep=3pt]
        \item\small (Palestra) GetNinjas: Da ideia a R\$47MM em investimento
        \item\small (Palestra) InternetcomResponsa: Cuidados e Responsabilidades no uso da Internet
        \item\small (Palestra) O papel da Inteligência Artificial na transformação dos negócios
        \item\small (Palestra) Tecnologia \& Empreendedorismo
        \item\small (Evento) 26º Encontro de Alunos e Ex-alunos da FIPP
    \end{itemize}


\end{itemize}    

\end{document}